\section{Preliminaries}\label{sec:prelim}

\paragraph{Setup and notation.}
We study an episodic, finite-horizon Markov decision process
$\mathrm{MDP}(\St,\Ac,H,\Pop,r)$ with state space $\St$, action space $\Ac$,
horizon $H$, transition kernels $\Pop=\{\Pop_h\}_{h\in[H]}$, and reward
functions $r=\{r_h\}_{h\in[H]}$ with $r_h(x,a)\in[0,1]$. An agent interacts for
$K$ episodes; we write $T:=KH$ for the total number of steps and
$[n]:=\{1,\dots,n\}$. For a policy $\pi=\{\pi_h\}_{h\in[H]}$, the value and
action-value functions are
\begin{align*}
V_h^\pi(x) \;=\; \E\!\Big[\textstyle\sum_{h'=h}^{H} r_{h'}(x_{h'},a_{h'}) \,\Big|\, x_h=x\Big],
\qquad
Q_h^\pi(x,a) \;=\; r_h(x,a) + (\Pop_h V_{h+1}^\pi)(x,a),
\end{align*}
where $(\Pop_h g)(x,a):=\E_{x'\sim \Pop_h(\cdot|x,a)}[g(x')]$ and
$V_{H+1}^\pi\equiv 0$. We write $V_h^\star,Q_h^\star$ for the optimal value
functions and $\T_h g(x,a):=r_h(x,a)+(\Pop_h g)(x,a)$ for the one-step Bellman
backup of a function $g:\St\to\R$. Throughout, $c, C, C_1, C_2, \ldots$ denote
universal positive constants whose values may change from line to line and
depend only on quantities specified in each statement's hypotheses; any
constant with a problem-dependent meaning is written with an explicit subscript
(e.g.\ $C_\beta$).

\paragraph{Regret.}
At episode $k$ the agent plays policy $\pi^k$ starting from a state $x_1^k$. The
cumulative regret over $K$ episodes is
\begin{align}\label{eq:regret-def}
\Reg(K) \;:=\; \sum_{k=1}^{K}\big[\, V_1^\star(x_1^k) - V_1^{\pi^k}(x_1^k) \,\big].
\end{align}

\begin{assumption}[Linear MDP]\label{ass:linear-mdp}
There exists a known feature map $\featmap:\St\times\Ac\to\R^d$ such that for
every $h\in[H]$ there exist $d$ unknown signed measures
$\mu_h=(\mu_h^{(1)},\dots,\mu_h^{(d)})$ over $\St$ and an unknown vector
$\theta_h\in\R^d$ with
\begin{align*}
\Pop_h(\cdot\,|\,x,a) \;=\; \inner{\featmap(x,a)}{\mu_h(\cdot)},
\qquad
r_h(x,a) \;=\; \inner{\featmap(x,a)}{\theta_h},
\qquad \forall (x,a)\in\St\times\Ac.
\end{align*}
We normalize $\norm{\featmap(x,a)}\le 1$ for all $(x,a)$ and
$\max\{\norm{\mu_h(\St)},\norm{\theta_h}\}\le\sqrt d$ for all $h\in[H]$, where
$\mu_h(\St):=\int_{\St} d\mu_h(x')\in\R^d$.
\end{assumption}

\begin{remark}\label{rem:linear-mdp-mild}
\Cref{ass:linear-mdp} is the standard linear-MDP model of
\cite{jin2020provably}. It strictly generalizes the tabular MDP (take
$\featmap(x,a)=e_{(x,a)}$, $d=|\St||\Ac|$) and the simplex feature model, yet
the transition kernel retains infinitely many degrees of freedom because
$\mu_h$ is unknown. The normalization $\norm{\featmap}\le1$,
$\norm{\theta_h},\norm{\mu_h(\St)}\le\sqrt d$ is a scaling convention rather
than a restriction: any bounded feature map can be rescaled to meet it.
\end{remark}

\paragraph{The LSVI-UCB algorithm.}
We analyze Least-Squares Value Iteration with an Upper-Confidence-Bound bonus
(\cite{jin2020provably}). Fix a regularizer $\lambda>0$ and a bonus radius
$\beta>0$. At episode $k$, for $h=H,H-1,\dots,1$ the algorithm forms the Gram
matrix and ridge weights
\begin{align}
\Gram_h^k &\;=\; \sum_{\tau=1}^{k-1} \featmap(x_h^\tau,a_h^\tau)\featmap(x_h^\tau,a_h^\tau)^\top + \lambda I, \label{eq:gram-def}\\
w_h^k &\;=\; (\Gram_h^k)^{-1} \sum_{\tau=1}^{k-1} \featmap(x_h^\tau,a_h^\tau)\big[\, r_h(x_h^\tau,a_h^\tau) + V_{h+1}^k(x_{h+1}^\tau) \,\big], \label{eq:weight-def}
\end{align}
and sets the optimistic action-value and value functions
\begin{align}
Q_h^k(x,a) &\;=\; \min\!\Big\{\, \inner{w_h^k}{\featmap(x,a)} + \beta\,\Mnorm{\featmap(x,a)}{(\Gram_h^k)^{-1}},\; H \,\Big\}, \label{eq:Q-def}\\
V_h^k(x) &\;=\; \max_{a\in\Ac} Q_h^k(x,a), \label{eq:V-def}
\end{align}
with $V_{H+1}^k\equiv 0$, where $\Mnorm{u}{M}:=(u^\top M u)^{1/2}$ is the
$M$-weighted (Mahalanobis) norm. The agent then plays the greedy policy
$\pi^k$, $a_h^k=\argmax_{a} Q_h^k(x_h^k,a)$. Let $\iota:=\log(2dT/\delta)$ for
the target failure probability $\delta\in(0,1)$; the analysis uses $\lambda=1$
and $\beta=C_\beta\,dH\sqrt{\iota}$ for a universal constant $C_\beta$ fixed in
\Cref{lem:concentration}.

\begin{remark}\label{rem:filtration}
We work on the filtration $\{\mathcal F_{k,h}\}$ generated by all transitions
observed up to and including step $h$ of episode $k$. The pair
$(\Gram_h^k,w_h^k)$ is $\mathcal F_{k-1,H}$-measurable: it depends only on the
first $k-1$ episodes. Consequently $V_h^k$ is determined before the fresh
transition at episode $k$ is revealed, a fact used in the martingale argument
of \Cref{sec:main}.
\end{remark}

The following structural fact is the foundation of the analysis: in a linear
MDP the Bellman backup of \emph{any} bounded function is linear in the features.

\begin{fact}[Linearity of the Bellman backup]\label{fac:value-linear}
Under \Cref{ass:linear-mdp}, for every $h\in[H]$ and every measurable
$g:\St\to[0,H]$ there is a vector $w\in\R^d$ with
\begin{align}\label{eq:backup-linear}
\T_h g(x,a) \;=\; \inner{\featmap(x,a)}{w},
\qquad
w \;=\; \theta_h + \int_{\St} g(x')\,d\mu_h(x'),
\qquad
\norm{w}\le 2H\sqrt d .
\end{align}
In particular, for any policy $\pi$ there exist weights $w_h^\pi$ with
$Q_h^\pi(x,a)=\inner{\featmap(x,a)}{w_h^\pi}$ and $\norm{w_h^\pi}\le 2H\sqrt d$.
\end{fact}

\begin{proof}
We compute the backup directly and then bound its weight vector. For the
identity,
\begin{align*}
\T_h g(x,a)
&\;=\; r_h(x,a) + \int_{\St} g(x')\,\Pop_h(dx'\,|\,x,a) \\
&\;=\; \inner{\featmap(x,a)}{\theta_h} + \int_{\St} g(x')\,\inner{\featmap(x,a)}{d\mu_h(x')} \\
&\;=\; \inner{\featmap(x,a)}{\theta_h} + \Inner{\featmap(x,a)}{\int_{\St} g(x')\,d\mu_h(x')} \\
&\;=\; \Inner{\featmap(x,a)}{\theta_h + \int_{\St} g(x')\,d\mu_h(x')},
\end{align*}
where the first step is the definition of $\T_h$, the second substitutes
\Cref{ass:linear-mdp} for both $r_h$ and $\Pop_h$, the third pulls the
$x'$-independent vector $\featmap(x,a)$ out of the integral by linearity of the
inner product, and the last collects terms. This proves
Eq.~\eqref{eq:backup-linear} with $w=\theta_h+\int g\,d\mu_h$. For the norm,
\begin{align*}
\norm{w}
&\;\le\; \norm{\theta_h} + \Norm{\int_{\St} g(x')\,d\mu_h(x')}
\;\le\; \sqrt d + H\,\norm{\mu_h(\St)}
\;\le\; \sqrt d + H\sqrt d
\;\le\; 2H\sqrt d,
\end{align*}
where the first step is the triangle inequality, the second uses
$0\le g\le H$ together with $\int d\mu_h=\mu_h(\St)$, the third applies the
normalization $\norm{\theta_h},\norm{\mu_h(\St)}\le\sqrt d$ from
\Cref{ass:linear-mdp}, and the last uses $H\ge1$. Applying the identity with
$g=V_{h+1}^\pi\in[0,H]$ and recalling $Q_h^\pi=\T_h V_{h+1}^\pi$ yields the
final claim. This completes the proof.
\end{proof}
